% ============================================================
\section{Discussion}
\label{sec:discussion}
% ============================================================

\paragraph{The operational failure mode is reproducible incorrectness.}
At baseline, no model in the fleet is reproducibly correct. Thirteen of seventeen
are reproducibly \emph{incorrect}: ten identical wrong answers out of ten, at
temperature $1.0$. This is the state that matters for a deployed system, and it is
worse than instability. A model that answers inconsistently is visible --- it fails
a repeated-execution check, and the system around it can be built to expect that.
A model that returns the same wrong action ten times in a row passes every such
check. The behaviour is stable, the test suite is green, and the specification
defect is invisible until the action is taken in production. Framing this task as
an accuracy problem understates it: the issue is not that models are wrong at some
rate, but that their wrongness is reproducible, and reproducibility is the property
systems are built to trust.

\paragraph{Instruction and specification are different levers, and only one moves the fleet.}
The seven conventional interventions are the standard toolkit: reasoning
elicitation, role assignment, encouragement, urgency, threat, error avoidance,
anti-hallucination. Applied to the same task, they move the fleet-wide rate from
$7.6\%$ to at most $36.5\%$. More tellingly, across all seven conditions and all
seventeen models, exactly one reproducibly incorrect model is ever converted to
reproducibly correct, and it is the same model each time. Expert framing
additionally stabilises a model that was non-reproducible at baseline, which is a
different and lesser achievement. The specification moves the same fleet to
$95.9\%$, with fourteen of seventeen models unanimous, eleven of them converted
directly from reproducible incorrectness. Every model moves toward the correct
action; none moves away.

The distinction is not that one prompt is longer or better written. Conventional
interventions address the model's disposition --- how carefully it should think, who
it should imagine itself to be, what it should avoid. They leave the task's
underdetermination untouched. The specification addresses the task: it states which
entities are involved, what condition satisfies the objective, and which actions are
feasible. The task is underdetermined, so the intervention that resolves the
underdetermination is the one that works.

\paragraph{What the internal measurements contribute, and what they do not.}
Section~\ref{sec:internal-measurement} is evidence that the behavioural change
corresponds to a change in the computation rather than to output formatting,
sampling variation, or a token-level coincidence at the answer position. The
substrate-conditioned state differs from the untreated state, and transplanting it
into the untreated computation reverses that computation's final preference. That is
the whole of the claim. These measurements do not identify a mechanism, do not
establish which constraints are internally represented, and are made on open-weight
models only --- the closed models in the behavioural fleet cannot be assumed to
realise the effect the same way.

\paragraph{The useful property is that the fix is a file.}
The intervention is a text artifact. It can be version-controlled, hashed,
diffed and reviewed like any other configuration; it applies to models whose weights
are inaccessible; it requires no training run, no labelled data and no evaluation
loop to fit. In this fleet it transferred across six vendors without per-model
adjustment. The cost is authoring: the constraints have to be identified and written
down correctly, which is human work and is not automated here. Fine-tuning can also
repair a task like this one, but it produces a new artifact per model, requires
weight access, and must be redone when the underlying model is replaced. For a
defect that originates in an underspecified input, specifying the input is the
proportionate repair.

\subsection{Limitations}

\paragraph{One task, one decision.}
All results are for a single binary operational decision with one correct answer.
We do not show that specification repairs multi-step tasks, tasks with several
interacting constraints, or tasks whose correct action is contested. The task was
chosen because it is small enough to verify exhaustively, not because it is
representative.

\paragraph{The specification is task-instantiated.}
The substrate names the entities of this task. It is not a general-purpose prefix,
and we make no claim that a specification written for one task helps on another.
Whether the effect survives abstraction is a separate question and is untested here.

\paragraph{Reproducibility is evaluated at $R=10$.}
Ten samples per model per condition establishes unanimity over ten draws, not over
the model's output distribution. A model recorded as reproducibly correct here could
produce a different action on the eleventh sample. The criterion is a practical bar,
not a probabilistic guarantee.

\paragraph{Closed models expose only their outputs.}
For the closed-weight majority of the fleet we observe the emitted text. Margins are
available only for the open-weight models, so the fleet-wide results and the
internal results are measured on overlapping but different populations, by different
instruments.

\paragraph{Three models were not repaired.}
The substrate leaves DeepSeek V3.2 at $8/10$, Llama 3.1-8B at $7/10$ and
Llama 3.2-3B at $8/10$. Two of these were reproducibly incorrect at baseline and
became non-reproducible rather than correct. The intervention improves every model
in the fleet and makes most of them reproducible; it does not make all of them
reproducible.

\paragraph{Routing.}
Fleet executions were issued through a single routing provider, so served snapshots
and backends are recorded but not controlled. Served model identifiers are reported
in Appendix~\ref{app:models}.

\subsection{Next}

The measurements reported here establish that the transition exists, that ordinary
prompting does not produce it, and that it has an internal counterpart. They do not
explain it. Which constraints carry the effect, whether the same constraints carry it
across model families, and how it depends on specification structure require
ablations across many specifications and models, together with attribution and
pathway analysis. That is the subject of separate work.
